\documentclass[a4paper]{article}
\usepackage{amsfonts}
\usepackage{amsmath}
\usepackage{amssymb}
\usepackage{graphicx}

\begin{document}
  
\noindent \large Auteur: Michael Livchits \\
\noindent \large Studierichting: Informatica\\
\noindent \large Oefeningenreeks 7E nr. 2.5\\

\medskip

\normalsize

$f(x,y)= (x^2+y^2)e^{-(x^2+y^2)}$\\
	
Eerste orde:\\
	
$\dfrac{\partial f}{\partial x} = 2xe^{-(x^2+y^2)(1-x^2-y^2)}$\\
	
$\dfrac{\partial f}{\partial y} = 2ye^{-(x^2+y^2)(1-x^2-y^2)}$\\

Kritiek punt: $(0,0)$ en alle punten in $x^2+y^2 = 1$\\

$f(x,y)= (x^2+y^2)e^{-(x^2+y^2)} \Rightarrow$ beide termen zijn altijd positief dus in $f(0,0)$ bereikt het een minimum.\\

Stel $x^2 + y^2 = m$\\

$f(m) = me^{-m}$\\

$f'(m) = e^{-m}(1-m)$\\

Nulpunt: $m = 1$\\

\begin{tabular}{c|ccc}
$m$     & & $1$ &\\ \hline
$f'(m)$ &$\nearrow$& $0$ &$\searrow$ \\
\end{tabular}\\

Alle punten in $x^2+y^2 = 1$ bereiken een maximum.\\

\begin{thebibliography}{999}
%\bibitem{a} Referentie
\end{thebibliography}
\end{document} 